\begin{table}[t]
\centering\fontsize{7}{8.5}\selectfont\setlength{\tabcolsep}{2.3pt}
\caption{\textbf{Comparing the library, its combinations, and additional baselines.}
All columns report mean relative $L^2$ error in percent on the same evaluation
cases. The best additional baseline is identified by its B label, and M1 to M9
are defined in Table~\ref{tab:paperbaselines}. The selected model uses one member,
the inverse error mixture weights individual errors, and the fitted mixture
optimizes their combination. All three use five fitting examples. Bold marks the lowest unrounded evaluation error.
Daggers mark exclusions in the separate sensitivity analysis
(Appendix~\ref{app:subsets}). Section~\ref{sec:experiments} specifies coverage.
Imported data: Poisson II, Heat II, and Navier Stokes \citep{niu2024}, and ERA5
\citep{hersbach2020,niu2024}.}
\label{tab:allexperts}
\resizebox{\linewidth}{!}{\begin{tabular}{lr|rrrrrrrrr|rrr}
\toprule
Dataset & Best baseline & M1 & M2 & M3 & M4 & M5 & M6 & M7 & M8 & M9 & \shortstack[r]{Selected\\model} & \shortstack[r]{Inverse\\error\\mixture} & \shortstack[r]{Fitted\\mixture} \\
\midrule
Darcy & 3.37 (B1) & 3.07 & 3.44 & 2.78 & 4.98 & \textbf{2.32} & 31.7 & 3.37 & 3.41 & 3.41 & 2.52 & 2.55 & 2.33 \\
Eikonal & 2.21 (B8) & 2.42 & 2.15 & 1.81 & 4.79 & 2.59 & 7.37 & 3.63 & 3.74 & 2.95 & 1.81 & 1.76 & \textbf{1.76} \\
Helmholtz & 4.12 (B9) & \textbf{3.3} & 3.48 & 4.79 & 10.9 & 5.04 & 29.1 & 7.1 & 10.9 & 10.8 & 6.25 & 6.89 & 6.55 \\
Poisson I & \textbf{0.00531} (B2) & 0.0413 & 0.0387 & 0.0427 & 0.0611 & 0.0955 & 9.92 & 0.00684 & 0.128 & 0.0838 & 0.00684 & 0.008 & 0.00788 \\
Poisson II$^{\dagger}$ & 0.0128 (B1) & 0.256 & 0.536 & 0.2 & 25.5 & 0.881 & 26.6 & \textbf{0.0121} & 0.921 & 0.798 & \textbf{0.0121} & 0.0125 & 0.0129 \\
Pressure Poisson$^{\dagger}$ & 0.0639 (B9) & 0.0599 & 0.0695 & 0.0918 & 0.0871 & \textbf{0.0498} & 0.285 & 0.0922 & 0.309 & 0.145 & 0.0734 & 0.0769 & 0.0802 \\
Heat I & 0.0206 (B10) & 0.0164 & 0.0125 & 0.0105 & 0.00915 & 0.0268 & 2.28 & 0.327 & 0.0494 & 0.0149 & 0.00972 & \textbf{0.00593} & 0.0062 \\
Heat II & 0.0201 (B10) & 0.0139 & 0.00878 & 0.00859 & 0.00662 & 0.0173 & 1.58 & 0.312 & 0.0268 & 0.0121 & 0.00662 & 0.00437 & \textbf{0.00435} \\
Cavity & \textbf{0.0419} (B5) & 98.5 & 0.29 & 98.5 & 0.681 & 3.6 & 70.1 & 0.409 & 2.06 & 0.164 & 0.164 & 0.142 & 0.142 \\
Navier Stokes & 1.52 (B10) & 1.55 & 1.69 & 1.57 & 1.58 & 2.33 & 4.41 & 2.79 & 2.62 & 1.95 & 1.61 & 1.52 & \textbf{1.5} \\
Rayleigh Bénard & 0.011 (B10) & 0.0112 & 0.00915 & 0.00715 & 0.00301 & 0.0137 & 0.0999 & 0.295 & 0.0529 & 0.0195 & 0.00301 & 0.00286 & \textbf{0.00269} \\
Burgers & 0.562 (B9) & 0.547 & 0.533 & 100 & 0.507 & 1.28 & 4.88 & 2.5 & 1.69 & 0.545 & 0.516 & 0.513 & \textbf{0.502} \\
Euler & \textbf{5.54} (B2) & 6.2 & 6.08 & 6.72 & 6.3 & 7.43 & 8.13 & 6.15 & 6 & 5.88 & 6.09 & 5.64 & 5.73 \\
Porous medium & 0.0718 (B9) & 0.0448 & 0.0325 & 0.0596 & 0.0654 & 0.126 & 2.75 & 0.909 & 0.681 & 0.0497 & 0.0325 & \textbf{0.0262} & 0.0263 \\
Allen Cahn & 41.2 (B5) & 53 & 52.7 & 42.1 & 45.4 & 41.3 & 102 & 55.5 & 41.4 & 38.8 & 39.4 & 37.9 & \textbf{34.6} \\
Cahn Hilliard I$^{\dagger}$ & 13.3 (B10) & 12.9 & 12.5 & 12.2 & 11.7 & 36.3 & 81.6 & 30.3 & 41.1 & 12.5 & 12.2 & 11.3 & \textbf{10.7} \\
Cahn Hilliard II & 49.8 (B5) & \textbf{46.9} & 48.3 & 51.9 & 50.8 & 60.8 & 99.8 & 64.8 & 55.9 & 47.8 & 52.2 & 47.2 & 49.3 \\
Fisher KPP & 3.24 (B9) & 3.21 & 3.22 & 2.36 & 2.98 & 3.83 & 6.49 & 6.86 & 2.86 & 2.84 & 2.36 & 2.53 & \textbf{2.3} \\
Phase field crystal & 83.5 (B11) & 89.6 & 89.5 & 111 & 83.3 & 202 & 120 & 83.5 & 85.3 & 85.7 & 83.5 & 83.8 & \textbf{83.3} \\
Shallow water & 0.403 (B10) & 0.428 & 0.367 & 0.362 & 0.876 & 0.446 & 1.38 & 1.09 & 1.29 & 0.502 & 0.367 & \textbf{0.357} & 0.389 \\
Wave & \textbf{3.3} (B3) & 6.36 & 6.31 & 4.75 & 11.1 & 5.04 & 64.2 & 24.9 & 8.35 & 11.3 & 4.84 & 4.39 & 4.27 \\
\midrule ERA5 & \textbf{5} (B1) & 14.2 & 17.7 & 20.4 & 12.2 & 9.75 & 9.54 & 6.74 & 20.4 & 8.24 & 6.74 & 5.81 & 6.49 \\
\bottomrule
\end{tabular}
}
\end{table}

\begin{table}[t]
\centering\fontsize{7.5}{9}\selectfont\setlength{\tabcolsep}{3pt}
\caption{\textbf{Model sources and performance, ranked by Elo.}
For each dataset, ratios divide a method's mean evaluation error by the
Selected model's error before geometric averaging. The Selected model is the
member of M1 to M9 with the lowest mean relative $L^2$ error on the five fitting
examples. \textbf{PDE class ratio} averages within each class
and then equally across the seven PDE classes (21 datasets).
\textbf{Dataset ratio} weights all 22 datasets equally, including ERA5.
A ratio of 1 matches selection, 0.90 means 10\% lower aggregate error, and
1.10 means 10\% higher aggregate error. Elo compares each pair of methods on all
22 datasets, with higher ratings better (Appendix~\ref{app:elo}).
M1 to M9 form the library, and B1 to B11
are additional baselines.
Citations identify components or ideas, with adaptations in
Appendix~\ref{app:adaptations}. Individual errors appear in
Table~\ref{tab:allexperts} and Appendix~\ref{app:matchedbaselines}.
Training costs and fine data budgets differ.}
\label{tab:paperbaselines}
\resizebox{\linewidth}{!}{\begin{tabular}{llrrr}
\toprule
Model / rule & Component or idea source & PDE class ratio & Dataset ratio & Elo \\
\midrule
Fitted mixture & \citep{breiman1996} & 0.884 & 0.925 & 1988 \\
Inverse error mixture & \citep{goel2007} & 0.888 & 0.932 & 1951 \\
Selected model & Fitting set selection & 1.000 & 1.000 & 1851 \\
M2 All pairs FNO & \citep{li2021fno,herde2024} & 1.333 & 1.541 & 1688 \\
M1 FiLM FNO transfer & \citep{li2021fno,perez2018} & 1.940 & 2.035 & 1642 \\
M9 ConvNeXt corrector & \citep{liu2022convnext} & 1.695 & 1.867 & 1634 \\
M3 ConvNeXt transfer & \citep{liu2022convnext} & 2.607 & 2.445 & 1619 \\
M4 Distribution FNO & \citep{yu2026fire} & 1.795 & 2.155 & 1592 \\
B10 FNO transfer II & \citep{lyu2023} & 1.794 & 1.850 & 1590 \\
B9 FNO transfer I & \citep{lyu2023} & 1.794 & 1.848 & 1578 \\
M5 Wavelet transfer & \citep{tripura2022wno} & 2.459 & 2.400 & 1528 \\
B5 Latent MF network & \citep{li2022dmfal} & 3.794 & 2.772 & 1497 \\
M7 POD GP & \citep{higdon2008} & 5.492 & 3.108 & 1450 \\
B1 Autoregressive POD GP & \citep{kennedy2000} & 5.854 & 3.362 & 1449 \\
M8 Slice attention corrector & \citep{wu2024transolver} & 4.189 & 3.549 & 1439 \\
B2 Nonlinear POD GP & \citep{perdikaris2017} & 5.871 & 3.626 & 1424 \\
B3 Composite MF MLP & \citep{meng2020} & 4.174 & 3.974 & 1387 \\
B8 Fidelity basis FNO & \citep{li2022ifc} & 8.979 & 6.217 & 1364 \\
B4 Composite MF DeepONet & \citep{howard2023} & 8.001 & 5.803 & 1337 \\
B11 Affine POD GP & \citep{zhang2018} & 17.624 & 11.236 & 1218 \\
B6 Nonlinear decoder transfer & \citep{seidman2022} & 9.927 & 8.252 & 1197 \\
M6 Retrieved field DeepONet & \citep{lu2022mf} & 20.847 & 14.890 & 1057 \\
B7 MFRNP adapter & \citep{niu2024} & 49.969 & 24.194 & 1020 \\
\bottomrule
\end{tabular}
}
\end{table}
